% Generated by scripts/generate-worked-calculations.R; do not edit by hand.
% Registry ID: hyp.mus.sample_size
% Source notebook: notebooks/support/hypothesis-testing/support.Rmd
In Table \ref{tab:attribute_sample_sizes_alpha05}, the sampling interval $SI_k$ for $k$ tolerated errors is calculated by dividing the total book value $X$ by the minimum required sample size $n_k$. For example, $SI_2 =$ 13,500,000 / 187 = 72,192.51. This implies that if we tolerate two errors and we find two errors, the point estimate equals $k \cdot SI_2 =$ 2 $\cdot$ 72,192.51 = 144,385.03. Tolerating two errors, therefore, is synonymous with building error tolerance for an expected error of 144,385.03.

Similarly, if the expected error is 95,744.68, the sample size $n =$ 141 and the critical region $\{k | k \leq 1\}$.

In \emph{Case: Accounts receivable circularization}, the expected error $EE =$ 100,000. Therefore, we need to find a $t^*$ between the sample sizes $n_1 =$ 141 and $n_2 =$ 187 such that the point estimate $t^* X / n^* = EE$.

Applying formula \ref{eq:MUS_sample_size} to the parameters from the \emph{Case: Accounts receivable circularization}, we obtain:

\begin{equation*}
n^* = \frac{1 - \frac{141}{187 - 141}}{\frac{100,000}{13,500,000} - \frac{1}{187 - 141}} = 144.1
\end{equation*}
Rounding this sample size upward leads to a required sample size $n =$ 145.\footnote{This calculation assumes that the sample is evaluated with the Stringer bound method. If the cell evaluation method is used, the sample size is usually lower, because errors with a small tainting have no effect on the precision achieved.} We can calculate the number of tolerable taintings from

\begin{equation*}
t_{tol} = k + \frac{EE - M_{U, k}}{M_{U, k+1} - M_{U, k}} = 1 + \frac{450,000 - 436,007}{577,535 - 436,007} = 1.09890
\end{equation*}
where $k$ and $k+1$ are the two numbers of errors between which the sample size has been interpolated.

The conclusion is that, if we expect a total error in the population of 100,000, the minimum sample size is 145, and we can tolerate one 100\% error and a partial error of 9.890\%.
